% ============================================================
%  Preámbulo: página, fuentes, colores, títulos y tablas
% ============================================================

% --- Página (carta, márgenes de 1 in como el original de Word) ---
\usepackage[letterpaper,margin=2.54cm]{geometry}

% --- Idioma y fuentes ---
\usepackage{fontspec}
\usepackage[spanish,es-noshorthands]{babel}
% Calibri si está instalada (Windows); si no, Carlito (métricamente idéntica).
% Sin ligaduras comunes (fi, ti, ...), igual que Word.
\IfFontExistsTF{Calibri}
  {\setmainfont{Calibri}[Ligatures={TeX,NoCommon}]}
  {\setmainfont{Carlito}[Ligatures={TeX,NoCommon}]}
\linespread{1.05}

% --- Paquetes ---
\usepackage[table]{xcolor}
\usepackage{array}
\usepackage{xltabular}
\usepackage{titlesec}
\usepackage{enumitem}
\usepackage{graphicx}
\usepackage{tikz}
\usetikzlibrary{arrows.meta,shadows}
\usepackage[hidelinks]{hyperref}
\hypersetup{pdftitle={H1 Análisis de contexto y drivers}}

% --- Colores (tema de Office) ---
\definecolor{titulo}{HTML}{1F3864}
\definecolor{subtitulo}{HTML}{2F5496}
\definecolor{encabezado}{HTML}{D9E2F3}

% --- Párrafos: alineados a la izquierda y sin división silábica (como Word) ---
\setlength{\parindent}{0pt}
\setlength{\parskip}{6pt}
\hyphenpenalty=10000
\AtBeginDocument{\raggedright}
\pagestyle{empty}

% --- Títulos: "1. Título" y "1.1 Título" ---
\titleformat{\section}
  {\color{titulo}\bfseries\fontsize{16}{19}\selectfont}{\thesection.}{0.3em}{}
\titleformat{\subsection}
  {\color{subtitulo}\bfseries\fontsize{13}{16}\selectfont}{\thesubsection}{0.3em}{}
\titlespacing*{\section}{0pt}{16pt}{6pt}
\titlespacing*{\subsection}{0pt}{12pt}{4pt}

% --- Listas con viñeta redonda ---
\setlist[itemize]{label={\small\textbullet}, leftmargin=1.2cm, itemsep=2pt, topsep=2pt}

% --- Tablas ---
\setlength{\tabcolsep}{4pt}
\setlength{\arrayrulewidth}{0.5pt}
\newcolumntype{L}[1]{>{\raggedright\arraybackslash}p{#1}}   % ancho fijo, texto a la izquierda
\newcolumntype{Y}{>{\raggedright\arraybackslash}X}          % ancho flexible, texto a la izquierda
\newcolumntype{C}{>{\centering\arraybackslash}X}            % ancho flexible, centrado (matrices)

% Celda de encabezado (negrita y centrada)
\newcommand{\cab}[1]{\centering\arraybackslash\bfseries #1}
% Fila de encabezado sombreada (el sombreado se recorta un poco a cada lado
% para que no tape las líneas verticales)
\newcommand{\filacab}{%
  \rowcolor{encabezado}[\dimexpr\tabcolsep-\arrayrulewidth\relax][\dimexpr\tabcolsep-\arrayrulewidth\relax]}

% Tabla con cortes de página y encabezado repetido
%   \begin{tabla}{|L{3cm}|Y|} ... \end{tabla}
\newenvironment{tabla}[1]
  {\par\footnotesize\renewcommand{\arraystretch}{1.3}%
   \setlength{\LTpre}{2pt}\setlength{\LTpost}{6pt}%
   \xltabular{\textwidth}{#1}}
  {\endxltabular\par}

% Variante para las matrices STK x CRN (letra más pequeña, filas más altas)
\newenvironment{matriz}[1]
  {\par\fontsize{7.5}{9}\selectfont\renewcommand{\arraystretch}{2}%
   \setlength{\LTpre}{2pt}\setlength{\LTpost}{6pt}%
   \xltabular{\textwidth}{#1}}
  {\endxltabular\par}
